Synthetic still images from modern generative models are increasingly difficult to separate from photographs by casual inspection, yet many content-moderation pipelines must run under tight latency and parameter budgets.
The practical question is therefore not only which architecture maximizes in-distribution (ID) accuracy, but which compact detector offers a usable accuracy--efficiency trade-off once images leave the training generators and pass through ordinary recompression.

Prior work has explored CNN detectors~\cite{wang2020cnn,tan2019efficientnet}, frequency-domain cues~\cite{frank2020freq}, and, more recently, state-space sequence models~\cite{gu2023mamba}.
Reported gains are often entangled with mismatched training protocols, ImageNet initialization, or heterogeneous evaluation sets, which complicates deployment decisions.
Frequency features are sometimes presented as broadly helpful for diffusion artifacts; likewise, strong ID scores can conceal fragile transfer to particular generators or content domains~\cite{ojha2023ufd}.

This paper contributes a \emph{protocol-locked empirical study} rather than a claim of state-of-the-art accuracy.
Under a single no-pretraining recipe, we compare:
\begin{enumerate}
\item compact CNN baselines (EfficientNet-B0, MobileNetV3-Small, ShuffleNetV2-x0.5)~\cite{tan2019efficientnet,howard2019mobilenetv3,ma2018shufflenet};
\item a gated frequency-aware CNN (LiteFreqNet~v2) and related ablations;
\item two study-specific pure-PyTorch SSM classifiers (\textbf{LiteSSM-A} and \textbf{LiteSSM-B}), with and without an attached frequency branch.
\end{enumerate}

Evaluation covers ID test performance, content-domain splits (CelebA-HQ vs.\ bedroom), UniversalFakeDetect (UFD) per-generator and Macro averages~\cite{ojha2023ufd}, JPEG quality-70 recompression, and measured efficiency (parameters, FLOPs, batch-1 latency, batch-32 throughput).
External pretrained reference detectors (UnivFD, NPR) are reported in a separate panel for contextualization only~\cite{ojha2023ufd,tan2024npr}.
An unmatched FLUX probe is deferred to the appendix and is not used for architecture claims~\cite{flux2024}.
We explicitly separate \textbf{UFD Macro AUC} from \textbf{OOD Pooled AUC} to avoid metric ambiguity.

\paragraph{Main finding.}
Under the evaluated latency budget, \textbf{LiteSSM-A} is the recommended compact operating point in the training-matched panel, improving average cross-generator behavior relative to the CNN suite, but \emph{no evaluated model}---including SSMs---removes generator-specific failures (notably DALL·E) or content-domain difficulty (bedroom).
Frequency enhancement is not universally beneficial under our protocol.

The remainder of the paper is organized as follows.
Section~\ref{sec:related} situates related detectors.
Section~\ref{sec:methods} details data, training, metrics, and reproducibility controls.
Section~\ref{sec:results} reports results.
Section~\ref{sec:discussion} discusses metric definitions, deployment implications, and boundaries.
Section~\ref{sec:conclusions} concludes.
